\chapter{Modern Processor Architectures}

\begin{khulasa}[: what will we learn in this chapter?]
\begin{itemize}
    \item How processors grew, and why the frequency stopped going up.
    \item Multi-core processors and parallel processing: the benefits and the problems, and Amdahl's law.
    \item Cache memory and the memory hierarchy, and why the order in which we read data matters so much.
\end{itemize}
\end{khulasa}

\begin{intabih}
    All the numbers in this chapter were \textbf{really measured} on a laptop with an Intel Core i7-8665U processor
    (4 cores, 8 threads, L1 cache = 32 KB and L2 = 256 KB per core, and a shared L3 = 8 MB). If you repeat the tests on
    your computer, your numbers will be a little different, but the effects will be the same.
\end{intabih}

%================================================
\section{Modern processors: trends and changes}
%================================================

\subsection{Fifty years of progress}
\begin{center}
\small
\begin{tabular}{|c|c|c|c|c|}
    \hline
    \ras{Processor} & \ras{Year} & \ras{Transistors} & \ras{Frequency} & \ras{Cores} \\ \hline
    Intel 4004 & 1971 & 2\,300 & 0.74 MHz & 1 \\ \hline
    Intel 8086 & 1978 & 29\,000 & 5--10 MHz & 1 \\ \hline
    Intel 80386 & 1985 & 275\,000 & 16--33 MHz & 1 \\ \hline
    Pentium 4 & 2000 & 42 million & 1.3--3.8 GHz & 1 \\ \hline
    Core i7-8665U & 2019 & about 2 billion & 1.9--4.8 GHz & 4 \\ \hline
    processors of 2024 & 2024 & tens of billions & about 5 GHz & from 8 to more than 100 \\ \hline
\end{tabular}
\end{center}

\subsection{The frequency wall}
Until about 2004, each year brought a faster processor because the \textbf{frequency} went up. Then the frequency
stopped at about 4 to 5 GHz. The main reason is \textbf{heat}:
\begin{itemize}
    \item Every clock pulse uses energy. The higher the frequency, the more energy is used, and energy turns into heat.
    \item At about 4 GHz the processor gets so hot that a normal computer cannot cool it.
\end{itemize}

\begin{tashbih}
    A runner runs faster and faster, until he reaches a speed he cannot keep without falling. Instead of asking him to
    run faster, we add other runners who share the distance with him. This is exactly what processor makers did.
\end{tashbih}

\subsection{Today's trends}
\begin{enumerate}
    \item \textbf{More cores} instead of a higher frequency (next section).
    \item \textbf{Bigger and smarter cache memory}, because memory became much slower than the processor (last section).
    \item \textbf{Special instructions}: SIMD (Chapter 4), and speed-up for artificial intelligence and encryption.
    \item \textbf{Different cores on the same chip}: strong cores for heavy work and economical cores for light work, to save battery.
\end{enumerate}

%================================================
\section{Multi-core processors and parallelism}
%================================================

\subsection{The ideas}
\begin{tarif}
    \begin{itemize}
        \item A \term{core}{cœur}: a complete processor (control unit, arithmetic unit, registers, L1 cache).
        \item A \term{multi-core processor}{processeur multi-cœur}: several cores on one chip.
        \item A \term{thread}{thread, fil d'exécution}: a series of instructions that runs on its own. A program may have one thread or several.
        \item \term{Parallel processing}{traitement parallèle}: running several parts of the work \textbf{at the same moment} on different cores.
    \end{itemize}
\end{tarif}

\begin{center}
\begin{tikzpicture}[core/.style={unit, minimum width=2.4cm, minimum height=1.5cm}]
    \draw[brandDark, very thick, rounded corners=8pt, fill=brandLight] (-6.1,-2.6) rectangle (6.1,2.1);
    \node[font=\bfseries] at (0,1.75) {a processor chip with 4 cores};
    \foreach \x/\i in {-4.5/1, -1.5/2, 1.5/3, 4.5/4} {
        \node[core] at (\x,0.6) {core \i\\{\scriptsize L1, L2}};
    }
    \node[unit=brandOrange, minimum width=11.2cm] at (0,-1.1) {shared cache memory L3};
    \node[unit=brandPurple, minimum width=4cm] (ram) at (0,-3.6) {main memory RAM};
    \draw[sahmT, <->] (0,-1.6) -- (ram);
\end{tikzpicture}
\end{center}

\begin{tashbih}
    A kitchen with one cook (one core) prepares one dish after another. A kitchen with 4 cooks (4 cores) prepares 4 dishes
    at the same time. But they share the same fridge (L3 and RAM), and one of them may wait for his turn in front of it.
    And one complex dish only gets faster if we can split it into independent tasks.
\end{tashbih}

\subsection{Hyper-Threading: one core for two threads}
\term{Simultaneous multithreading}{multithreading simultané} (SMT, or Hyper-Threading) makes one core look like two cores:
when a thread waits for data from memory, the second thread uses the free arithmetic units. This is why our processor shows
8 ``logical processors'' even though it has only 4 cores. But if the two threads need the same arithmetic unit all the time,
we gain almost nothing, as we will see in the test.

\subsection{A test: the same work on 1, 2, 4 and 8 threads}
The program below computes the sum of the square roots of 400 million numbers, and shares the work between $k$ threads:

\cfile{ch5/threads.c}
\screen{ch5/threads.out}

\begin{sharh}
    \s{18--26}{Each thread takes the numbers $id$, $id + k$, $id + 2k$\dots so it computes a part of the sum that does not overlap the others.}
    \s{34--35}{\code{pthread\_create}: start $k$ threads; each one runs on a core.}
    \s{36--39}{\code{pthread\_join}: wait for each thread to finish, then add the partial results.}
    \s{40--42}{Speed-up = time with one thread $\div$ time with $k$ threads.}
\end{sharh}

\begin{center}
\begin{tikzpicture}[x=1.1cm, y=1.1cm]
    \draw[->, thick] (0,0) -- (8.6,0) node[right, font=\small] {threads};
    \draw[->, thick] (0,0) -- (0,4.4) node[above, font=\small] {speed-up};
    \foreach \y in {1,2,3,4} \draw (0.08,\y) -- (-0.08,\y) node[left, font=\scriptsize] {\y};
    \foreach \x in {1,2,4,8} \draw (\x,0.08) -- (\x,-0.08) node[below, font=\scriptsize] {\x};
    \draw[brandMuted, dashed, thick] (1,1) -- (4,4) node[above, font=\scriptsize, text=brandMuted] {ideal speed-up};
    \draw[orangeDeep, very thick] (1,1) -- (2,1.98) -- (4,3.02) -- (8,3.05);
    \foreach \x/\y in {1/1, 2/1.98, 4/3.02, 8/3.05} {
        \fill[orangeDeep] (\x,\y) circle (2.5pt);
        \node[font=\codefont\scriptsize, below right] at (\x,\y) {\y};
    }
    \draw[brandTeal, dashed] (4,0) -- (4,4.2);
    \node[font=\scriptsize, text=tealDeep, align=center] at (6,3.9) {only 4 cores:\\extra threads do not help here};
\end{tikzpicture}
\end{center}

\begin{khatwa}[: how do we read these results?]
    \begin{enumerate}
        \item \textbf{2 threads}: speed-up 1.98, very close to 2. The work was split into two fully independent halves.
        \item \textbf{4 threads}: speed-up 3.02, not 4. When the four cores work together the processor gets hot, so it lowers
              its frequency by itself. And a part of the program (starting the threads and adding the results) is not parallel.
        \item \textbf{8 threads}: speed-up 3.05, almost the same as 4 threads. Our processor has only 4 real cores, and the
              extra threads share the same arithmetic units, so the work does not go faster.
    \end{enumerate}
\end{khatwa}

\subsection{Amdahl's law: the limits of parallelism}
Most programs have a part that \textbf{cannot be shared} (reading the file, showing the result, waiting for the user\dots).

\begin{tarif}
    If the parallel part of the program is $p$ (between 0 and 1), and we have $n$ cores, the maximum speed-up is:
    \[ S(n) = \frac{1}{(1 - p) + \dfrac{p}{n}} \]
\end{tarif}

\begin{khatwa}[: a program that is 90\% parallel]
    $p = 0.9$, and the serial part $1 - p = 0.1$.
    \begin{enumerate}
        \item With 4 cores: $S = \dfrac{1}{0.1 + \frac{0.9}{4}} = \dfrac{1}{0.1 + 0.225} = \dfrac{1}{0.325} \approx 3.08$.
        \item With 16 cores: $S = \dfrac{1}{0.1 + 0.05625} \approx 6.4$.
        \item With an infinite number of cores: $\dfrac{p}{n} \rightarrow 0$, so $S = \dfrac{1}{0.1} = 10$.
    \end{enumerate}
    However many cores we add, the program will never be more than 10 times faster, because the 10\% serial part stays the same.
\end{khatwa}

\begin{hal}
    Notice that Amdahl's law with $p = 0.9$ and 4 cores gives 3.08, very close to the speed-up of 3.02 we measured!
    A simple law explains reality well.
\end{hal}

\subsection{Benefits and challenges}
\begin{center}
\small
\begin{tabularx}{\linewidth}{|R|R|}
    \hline
    \ras{Benefits} & \ras{Challenges} \\ \hline
    more performance without raising the frequency & the program must be written with several threads, which is harder \\ \hline
    less energy for the same work & the serial part limits the speed-up (Amdahl's law) \\ \hline
    several programs run smoothly at the same time & \textbf{synchronization}: two threads writing the same variable can break the result \\ \hline
    if one task fails, the other cores keep working & \textbf{cache coherence}: all the cores must see the same value of a variable \\ \hline
\end{tabularx}
\end{center}

\begin{mithal}[: the synchronization problem]
    Two threads run \code{x = x + 1} on the same variable, and x = 5 at the start:
    \begin{enumerate}
        \item Thread 1 reads x = 5.
        \item Thread 2 also reads x = 5 (before the first one writes).
        \item Thread 1 computes $5 + 1$ and writes 6.
        \item Thread 2 computes $5 + 1$ and also writes 6.
    \end{enumerate}
    The result is 6 instead of 7: one increase is lost! This is why we use \textbf{locks} (mutex) or \textbf{atomic}
    instructions (\code{lock add} in x86) that stop two threads from changing the variable at the same time. In our program
    we avoided the problem: each thread writes into its own box \code{part[id]}.
\end{mithal}

%================================================
\section{Cache memory and the memory hierarchy}
%================================================

\subsection{The problem: memory is very slow}
The processor can run several instructions every nanosecond, but reading the main memory takes about 100 nanoseconds.
If the processor waited for memory every time, it would do nothing most of the time.

\begin{tarif}
    \term{Cache memory}{mémoire cache, antémémoire}: a small and very fast memory (SRAM) close to the processor. It keeps a
    \textbf{copy} of the data used recently, or of the data we will need soon.
\end{tarif}

\begin{tashbih}
    When you revise for an exam, you do not put all the books of the library on your desk. On the desk you put the books you
    need now (L1), in your bag the books of the week (L2), in your cupboard at home the books of the year (L3), and the rest
    stays in the far library (RAM). The closer the place, the smaller and the faster it is.
\end{tashbih}

\subsection{The memory hierarchy}
\begin{center}
\begin{tikzpicture}
    \foreach \y/\w/\t/\c/\d in {
        3.6/1.6/{registers}/brandOrange/{less than 1 KB, less than 1 ns},
        2.7/2.8/{L1}/brandOrange!70!brandTeal/{32 KB, about 1 ns},
        1.8/4.0/{L2}/brandTeal/{256 KB, about 4 ns},
        0.9/5.2/{L3}/brandTeal!70!brandPurple/{8 MB, about 10 ns},
        0.0/6.4/{main memory RAM}/brandPurple/{8 GB, about 100 ns},
        -0.9/7.6/{SSD disk}/brandMuted/{256 GB, about 100 µs}} {
        \fill[\c!35, draw=\c, thick] (-\w/2,\y) rectangle (\w/2,\y+0.8);
        \node[font=\small\bfseries] at (0,\y+0.4) {\t};
        \node[font=\scriptsize, anchor=west] at (4.1,\y+0.4) {\d};
    }
    \draw[sahmO] (-4.3,-0.9) -- (-4.3,4.4) node[midway, left, font=\small, align=right] {faster\\smaller\\more expensive};
\end{tikzpicture}
\end{center}

\subsection{A test: we can see the memory levels}
The program below reads an array in a random order (each read gives the address of the next read, so the processor cannot
guess), and measures the average time of one read, for different array sizes:

\cfile{ch5/latency.c}
\screen{ch5/latency.out}

\begin{center}
\begin{tikzpicture}[x=1.45cm, y=0.042cm]
    \draw[->, thick] (0,0) -- (7.4,0) node[right, font=\small] {data size};
    \draw[->, thick] (0,0) -- (0,125) node[above, font=\small] {time (ns)};
    \foreach \y in {0,25,50,75,100} \draw (0.06,\y) -- (-0.06,\y) node[left, font=\scriptsize] {\y};
    \foreach \x/\t in {1/16 KB, 2/64 KB, 3/256 KB, 4/1 MB, 5/4 MB, 6/16 MB, 7/64 MB}
        \node[font=\codefont\scriptsize, below] at (\x,-2) {\t};
    \fill[brandOrange!15] (0.5,0) rectangle (1.6,125); \node[font=\scriptsize] at (1.05,118) {L1};
    \fill[brandTeal!15] (1.6,0) rectangle (3.4,125); \node[font=\scriptsize] at (2.5,118) {L2};
    \fill[brandPurple!12] (3.4,0) rectangle (5.6,125); \node[font=\scriptsize] at (4.5,118) {L3};
    \fill[brandMuted!12] (5.6,0) rectangle (7.3,125); \node[font=\scriptsize] at (6.45,118) {RAM};
    \draw[orangeDeep, very thick] (1,1.7) -- (2,2.2) -- (3,4.1) -- (4,9.2) -- (5,12.8) -- (6,83.5) -- (7,112.7);
    \foreach \x/\y in {1/1.7, 2/2.2, 3/4.1, 4/9.2, 5/12.8, 6/83.5, 7/112.7} {
        \fill[orangeDeep] (\x,\y) circle (2.5pt);
        \node[font=\codefont\scriptsize, above] at (\x,\y+2) {\y};
    }
\end{tikzpicture}
\end{center}

\begin{khatwa}[: what do we see in the curve?]
    \begin{enumerate}
        \item \textbf{Up to 32 KB}: the whole array is in L1, and a read takes only 1.7 ns.
        \item \textbf{64 KB and 256 KB}: the array does not fit in L1 any more, so we read from L2: the time goes up to 2 then 4 ns.
        \item \textbf{1 MB and 4 MB}: we read from the shared L3: about 9 to 13 ns.
        \item \textbf{16 MB and more}: the array is bigger than L3 (8 MB), so we go to RAM: the time jumps to 84 then 113 ns,
              about \textbf{65 times} slower than L1!
    \end{enumerate}
    The steps in the curve fall exactly at the cache sizes of our processor.
\end{khatwa}

\subsection{The principle of locality: why does cache memory work?}
Cache memory is small, but it usually holds what we need, thanks to two natural habits of programs:
\begin{itemize}
    \item \term{Temporal locality}{localité temporelle}: the data we used now will often be used \textbf{again soon}
          (the loop counter, the sum\dots).
    \item \term{Spatial locality}{localité spatiale}: after reading an address, we often read the \textbf{nearby addresses}
          (array elements, instructions that follow each other).
\end{itemize}

\begin{tarif}
    A \term{cache line}{ligne de cache}: the cache does not move one byte, but a whole block of \textbf{64 bytes} next to
    each other. When we read \code{t[0]}, \code{t[1]} to \code{t[15]} (for 4-byte numbers) come with it for free.
\end{tarif}

\subsection{A test: the order of reading changes everything}
\cfile{ch5/locality.c}
\screen{ch5/locality.out}

\begin{center}
\begin{tikzpicture}[c/.style={draw=brandDark, minimum size=0.55cm}]
    \foreach \i in {0,...,3} \foreach \j in {0,...,3} \node[c] at (\j*0.55-6.0, -\i*0.55) {};
    \foreach \i in {0,...,3} \draw[sahmT] (-6.2, -\i*0.55) -- (-4.1, -\i*0.55);
    \node[font=\small, align=center] at (-5.2,0.75) {row by row: each read\\is next to the one before};
    \foreach \i in {0,...,3} \foreach \j in {0,...,3} \node[c] at (\j*0.55+3.0, -\i*0.55) {};
    \foreach \j in {0,...,3} \draw[sahmO] (\j*0.55+3.0, 0.2) -- (\j*0.55+3.0, -1.85);
    \node[font=\small, align=center] at (3.8,0.75) {column by column: each read\\jumps 16 KB};
\end{tikzpicture}
\end{center}

\begin{khatwa}[: why are columns 33 times slower?]
    \begin{enumerate}
        \item In C, the array \code{m[4096][4096]} is stored \textbf{row after row}: \code{m[0][0]}, then \code{m[0][1]}, then\dots
        \item \textbf{Row by row}: each cache line (64 bytes) holds 16 numbers next to each other. We read from RAM once, then use
              16 numbers from the cache. And the processor notices we move forward regularly, so it fetches the next lines in advance.
        \item \textbf{Column by column}: \code{m[0][j]} and \code{m[1][j]} are $4096 \times 4 = 16\,384$ bytes apart. Every read is
              in a new cache line, so we use only one number out of the 16 we moved.
        \item The real result: 0.009 s against 0.301 s, about \textbf{33 times} slower, for exactly the same operations!
    \end{enumerate}
\end{khatwa}

\subsection{Hits, misses and average access time}
\begin{tarif}
    \begin{itemize}
        \item A \term{hit}{succès}: the data is in the cache.
        \item A \term{miss}{défaut}: the data is not in the cache, so we go to the next (slower) level.
        \item The \textbf{hit rate} $h$: the share of reads that found their data in the cache.
    \end{itemize}
    \[ \text{average time} = h \times t_{cache} + (1 - h) \times t_{RAM} \]
\end{tarif}

\begin{khatwa}[: $t_{cache} = 1$ ns, $t_{RAM} = 100$ ns]
    \begin{enumerate}
        \item Hit rate 90\%: $0.9 \times 1 + 0.1 \times 100 = 0.9 + 10 = 10.9$ ns.
        \item Hit rate 99\%: $0.99 \times 1 + 0.01 \times 100 = 0.99 + 1 = 1.99$ ns.
    \end{enumerate}
    Going from 90\% to 99\% makes memory \textbf{more than 5 times faster}. This is why modern processors work hard to improve the hit rate.
\end{khatwa}

\subsection{Where do we put data in the cache? (mapping)}
The cache is split into lines. We must decide in which line we put a block coming from memory:

\begin{center}
\small
\begin{tabularx}{\linewidth}{|>{\centering\arraybackslash}p{3.4cm}|R|R|}
    \hline
    \ras{Kind} & \ras{Rule} & \ras{Advantage and drawback} \\ \hline
    direct mapped & each block has only one possible line: block number mod number of lines & simple and fast, but two blocks may fight for the same line \\ \hline
    fully associative & any block in any line & fewer conflicts, but searching all the lines is expensive \\ \hline
    set associative & the block goes to a given set, and inside it to any of $k$ lines & the middle way, used in all modern processors \\ \hline
\end{tabularx}
\end{center}

\begin{khatwa}[: splitting an address in a direct-mapped cache]
    A 1 KB cache, with 16-byte lines, and 16-bit addresses.
    \begin{enumerate}
        \item Number of lines: $1024 \div 16 = 64 = 2^6$ lines, so we need \textbf{6 bits} for the line number.
        \item Inside the line there are $16 = 2^4$ bytes, so \textbf{4 bits} for the offset.
        \item The rest: $16 - 6 - 4 = 6$ bits, called the \textbf{tag}: it tells which memory block is in the line now.
    \end{enumerate}
    \begin{center}
    \begin{tikzpicture}[f/.style={draw=brandDark, thick, minimum height=0.7cm, font=\small\bfseries}]
        \node[f, minimum width=3cm, fill=brandOrange!20] (t) at (0,0) {tag: 6 bits};
        \node[f, minimum width=3cm, fill=brandTeal!20, anchor=west] (l) at (t.east) {line: 6 bits};
        \node[f, minimum width=2cm, fill=brandGreen!20, anchor=west] (o) at (l.east) {offset: 4 bits};
        \node[font=\codefont\scriptsize, below=1pt of t.south west, anchor=north west] {bit 15};
        \node[font=\codefont\scriptsize, below=1pt of o.south east, anchor=north east] {bit 0};
    \end{tikzpicture}
    \end{center}
    Example: $\hex{1A2C} = \bin{000110\,100010\,1100}$. Tag $\bin{000110} = 6$, line $\bin{100010} = 34$, and offset $\bin{1100} = 12$.
\end{khatwa}

\subsection{Replacement and writing}
\begin{itemize}
    \item \textbf{Replacement policy}: when a set is full, which line do we throw out? Usually \textbf{LRU} (Least Recently Used):
          the line that has not been used for the longest time.
    \item \textbf{Write policy}:
    \begin{itemize}
        \item \textbf{Write-through}: we write in the cache and in RAM together. Simple but slow.
        \item \textbf{Write-back}: we write only in the cache, and write in RAM only when the line is thrown out. Faster, and used today.
    \end{itemize}
\end{itemize}

\begin{hal}
    Remember the Harvard model from Chapter 1: the L1 cache is split into two parts, L1i for instructions and L1d for data
    (32 KB each in our processor), so the processor can fetch an instruction and read data at the same time. But L2, L3 and RAM
    are shared by both, as in the Von Neumann model.
\end{hal}

%================================================
\section{Exercises}
%================================================

\begin{tamrin}
    A program is 75\% parallel. Compute the speed-up with 2, 4 and 8 cores, and the maximum speed-up.
\end{tamrin}
\begin{hall}
    $1 - p = 0.25$. \quad $S(2) = \dfrac{1}{0.25 + 0.375} = 1.6$, \quad $S(4) = \dfrac{1}{0.25 + 0.1875} \approx 2.29$,
    \quad $S(8) = \dfrac{1}{0.25 + 0.09375} \approx 2.91$. Maximum $= \dfrac{1}{0.25} = 4$.
\end{hall}

\begin{tamrin}
    A cache with a time of 2 ns, a main memory of 80 ns, and a hit rate of 95\%. What is the average time?
\end{tamrin}
\begin{hall}
    $0.95 \times 2 + 0.05 \times 80 = 1.9 + 4 = 5.9$ ns.
\end{hall}

\begin{tamrin}
    A direct-mapped cache: 4 KB, 32-byte lines, 32-bit addresses. How many bits for the offset, the line and the tag?
\end{tamrin}
\begin{hall}
    Offset: $32 = 2^5$, so 5 bits. Lines: $4096 \div 32 = 128 = 2^7$, so 7 bits. Tag: $32 - 7 - 5 = 20$ bits.
\end{hall}

\begin{tamrin}
    Explain why this code is slow, and how to fix it:
    \begin{codebox}\codefont\small
    for (j = 0; j < N; j++) \\
    \hspace*{2em}for (i = 0; i < N; i++) \\
    \hspace*{4em}s += m[i][j];
    \end{codebox}
\end{tamrin}
\begin{hall}
    It reads the array column by column, so every read is in a different cache line (no spatial locality). We fix it by
    swapping the two loops: \code{for (i ...) for (j ...)}, so we read row by row.
\end{hall}

\begin{tamrin}
    Why did the program \code{threads.c} not get better when going from 4 to 8 threads on our processor?
\end{tamrin}
\begin{hall}
    Because the processor has only 4 real cores. Threads 5 to 8 run with Hyper-Threading on the same cores, and the program
    uses the arithmetic unit all the time, so there are no free units for the second thread to use.
\end{hall}
